% TOP_PROBLEM: {"id": 3185, "title": "Seagull problem", "category_id": 3, "category": {"id": 3, "name": "graph_theory", "display_name": "Graph Theory", "description": "Problems involving graphs, networks, and their properties.", "slug": "graph-theory", "order_index": 3, "created_at": "2026-07-31T15:26:25.671Z"}, "status": "open", "problem_number": "OPG-729", "set_id": 12}
% FIRST_POSTED: 2026-10-01
% ENTRY_KIND: statement_only
% UnsolvedMath ID 3185; source catalogue code OPG-729
% !TeX program = lualatex
% Definition and sources only; this file is not an LLM solution attempt.
\documentclass[11pt,a4paper]{article}
\usepackage[margin=25mm]{geometry}
\usepackage{fontspec}
\setmainfont{lmroman10-regular.otf}[BoldFont=lmroman10-bold.otf,
 ItalicFont=lmroman10-italic.otf,BoldItalicFont=lmroman10-bolditalic.otf]
\usepackage{amsmath,amssymb,mathtools}
\usepackage{unicode-math}
\setmathfont{latinmodern-math.otf}
\AtBeginDocument{\let\setminus\smallsetminus}
\usepackage{xurl,hyperref}
\hypersetup{colorlinks=true,urlcolor=blue}
\setcounter{secnumdepth}{0}
\setlength{\parindent}{0pt}
\setlength{\parskip}{5pt}
\setlength{\emergencystretch}{3em}
\begin{document}
\section*{Seagull problem}\label{3185}
{\small UnsolvedMath ID 3185 · OPG-729 · Graph Theory · OpenGarden}
\subsection{Definitions and mathematical statement}
Let $G=(V,E)$ be a finite simple graph with
$n=|V|\geq1$. An independent set is a subset of vertices
with no edges between any two of them. The hypothesis
that there is no independent set of size three is
equivalent to saying that among every three distinct
vertices some pair is adjacent, or that the independence
number $\alpha(G)$ is at most two.

For an integer $t\geq1$, a $K_t$ minor can be specified
by pairwise disjoint nonempty branch sets
$B_1,\ldots,B_t\subseteq V$, such that each $G[B_i]$
is connected and there is an edge between every pair
of distinct branch sets. Contracting each branch set
and deleting surplus vertices and edges produces the
complete graph $K_t$.

The conjecture is
\[
 \alpha(G)\leq2\quad\Longrightarrow\quad
 G\text{ has a }K_{\lceil n/2\rceil}\text{ minor}.
\]
The ceiling gives the precise integer interpretation
of a complete minor on at least $n/2$ vertices.
No connectivity assumption is imposed on $G$, and
the branch sets need not exhaust its vertices.

The name refers to a seagull, an induced three-vertex
path: its two ends are nonadjacent and its middle
vertex is adjacent to both. Such connected triples
can serve as branch sets in attempts to build the
required minor, but the conjecture places no shape
or size restriction on its branch sets. It asks
for a complete minor, not just a dense minor or
a large clique subgraph.
\subsection{Short English statement}
Must every n-vertex graph with no three pairwise nonadjacent vertices have a complete minor on at least half its vertices?
\subsection{Sources}
\begin{itemize}
\item[S1] ULAM.AI and the UnsolvedMath Contributors, supplied problems.json, record 3185 (OPG-729), OpenGarden collection. Catalogue identity and source transcription; CC BY 4.0.
\url{https://huggingface.co/datasets/ulamai/UnsolvedMath}.
\item[S2] Open Problem Garden, Seagull problem. Original problem and discussion; the mathematical formulation below is expanded from the supplied source record.
\url{http://www.openproblemgarden.org/op/seagull_problem}.
\end{itemize}
\end{document}
